\subsection{Using verifier scores to admit or rank evidence (E6)}
\label{sec:e6}
\label{sec:inflation}\label{sec:eval-verify}\label{sec:exp-e6}
A score that is suitable for admitting passages need not work as a ranking
score. The hybrid tool of Section~\ref{sec:gate} uses the verifier score in
both roles, so we tested the second one directly. We took 16 candidate pools
from the historical hybrid wrapper and changed only their order. Under
verifier-score order, the number of pools with a gold document in the top
five fell from 15 to seven: eight pools lost a hit and none gained one
(exact McNemar $p=0.0078$; Figure~\ref{fig:verifier}b).

The admission question remains open. All 191 scores in a separate archive
exceeded the 0.5 cutoff (Figure~\ref{fig:verifier}a), so that archive says
little about the value of filtering. In a reconstructed replay with fixed
candidates, the filtered and unfiltered arms retained the same number of
documents on nine of ten questions; the archive does not establish that they
retained the same documents. The reordering result therefore does not show
whether filtering helps.

\subsection{What answer traces and success status reveal (E7)}
\label{sec:e5}
\label{sec:agentic}\label{sec:abstention}\label{sec:composition}\label{sec:eval-agent}
\label{sec:sysarms}\label{sec:e7}\label{sec:exp-e7}
On a shared five-collection corpus, each of three interfaces included
citation markers in nine of ten answers. A marker records an attempt to
cite; it does not show that the source supports the claim. Abstention was
not uniform. In three industrial-domain probes with no in-corpus answer, the
instructed agent reported not found in all three completed runs; this is an
observation of instructed agent behaviour, not a server-side guarantee. In a
separate probe family, the retrieval layer still returned candidates above
the 0.35 threshold for absent topics, so the threshold is not a test of
answerability.

The deployment probes exposed a different ambiguity: global search could
report success after skipping unreadable collections. Scoped search with
forced reindexing recovered a missing hit, and pagination recovered a shelf
listing cut off by the output limit. Reporting the scope searched and any
skipped collections makes it possible to distinguish incomplete search from
a completed search that found no evidence.

Both observations bound how the earlier results should be read. A target
string in the pack measures text coverage (Figure~\ref{fig:content}), and a
top-five hit measures document position (Figure~\ref{fig:verifier}); neither
establishes that a cited passage supports the answer.

% The verifier result remains a full-width float after E7 so that its
% caption stays adjacent to the paired score plot.
\begin{figure*}[!tb]
\centering
\includegraphics[width=0.92\textwidth]{fig_opt_verifier}
\caption{Verifier scores and same-pool ordering from separate archives.
(a) Laya scores for 191 recorded candidates across 16 hybrid runs. Median
scores were 0.88 for the 17 gold candidates and 0.86 for the 174 others.
All recorded scores exceeded 0.5; one more candidate (Q14) had no recorded
score. Horizontal jitter separates the marks. (b) In the separate ordering
archive, each pool is replayed in fusion and verifier-score order. Filled
marks indicate a gold document among the first five.}
\label{fig:verifier}
\end{figure*}
